% BOOK-FILE-ID: f-bcb250cd9e78
\section{行き詰まったとき}\label{節：行き詰まったとき}

式変形をしても一向に簡単にならない,ということはよくある.
しかし,そのとき残った項こそが次の手掛かりになる.例えば,
\begin{enumerate}
  \item \bookterm{logarithm}{対数}を取っても,\bookterm{sequence}{数列}の項を1次の和として表せない.
  \item 定数を引いて付加項を消そうとしたが,条件を満たす定数が求まらない.
  \item 置換後に新しい種類の項が増え続ける.
\end{enumerate}
定数で消せない$n$の項が残ったなら,1次式を引いてみる.対数を取っても和にならなければ,先に因数分解や定数の調整を試してみる.
残った項が何者なのかを見極めれば,次に打つ手が見えてくるのだ.
\BookFigFallback

行き詰まったときこそ,実験に立ち返ろう.\,$a_2$から$a_5$程度を計算すれば,一般項の予想が立つこともある.
ただし,予想は\bookterm{initial}{初期条件}と漸化式を使って,数学的帰納法などで証明しなければならない.
1つの方法がうまくいかなくても,その漸化式が解けないと決まったわけではない.別の公式や初項の値にも注目してみよう.

\bookterm{general}{一般項}が難しくても,正値・単調性・\bookterm{limit}{極限}・保存される量は調べられるかもしれない.
\bookproblemref{practice-entrance}{6}{発展問6}では逆数を取った漸化式から減り方を求める.
求めたい性質に合わせて,何を調べるのかを選び直すことも大切である.
